% 1Meter, DDS metrology, 1PDB, and gateway models — technical white paper
% Build: pdflatex (twice if you add \tableofcontents later)
\documentclass[11pt,a4paper]{article}
\usepackage[margin=2.2cm]{geometry}
\usepackage[T1]{fontenc}
\usepackage[utf8]{inputenc}
\usepackage{lmodern}
\usepackage{microtype}
\usepackage{hyperref}
\usepackage{booktabs}
\usepackage{tabularx}
\usepackage{parskip}
\usepackage{enumitem}

\hypersetup{colorlinks=true, linkcolor=blue, urlcolor=blue, citecolor=blue,
  pdftitle={1Meter, 1PDB, and gateway architecture},
  pdfauthor={1PWR Africa}}

\title{\textbf{1Meter, DDS Metrology, 1PDB,\\and Payment vs.\ Mesh Gateways}\\
\large Technical status, application architecture, and market impact}
\author{1PWR Customer Care / 1PDB program context}
\date{May 2026}

\begin{document}
\maketitle
\thispagestyle{empty}

\begin{abstract}
\noindent
This note summarizes the \textbf{1Meter} prototype (ESP32-C3, Mesh-Lite, AWS IoT), the \textbf{DDS8888} energy metrology path, the \textbf{1PDB} PostgreSQL backend behind Customer Care, and how \textbf{two different ``gateway'' concepts}---\emph{merchant payment gateway} versus \emph{mesh uplink gateway}---fit together. It reflects the architecture documented in the \texttt{1PWR~CC} repository; see \texttt{CONTEXT.md} and \texttt{docs/ops/1meter-billing-migration-protocol.md}.
\end{abstract}

\section{Purpose and strategy}

Rural minigrids need auditable kWh accounting at unit economics that scale. \textbf{SparkMeter} (Koios cloud and ThunderCloud on-prem) remains the regulated billing instrument across Lesotho and Benin. \textbf{1Meter} is 1PWR's in-house meter stack: owned firmware, cloud path, and hardware roadmap, aimed at lower hardware cost, full control of cutoff logic, and eventual SparkMeter-independent sites.

The \textbf{MAK} site in Lesotho is the controlled pilot: 1Meters installed in series with SparkMeters, with SparkMeter \textbf{primary} for billing until statistical and operational gates are met (Phase~1 of the billing migration protocol).

\section{Hardware: PCB and DDS8888}

The electrical design lives in the \textbf{1Meter\_PCB} KiCad repository. PCB revision \textbf{v2.2} connects the energy measurement IC to the ESP32 \textbf{only via RS485 Modbus}. The field stack uses a \textbf{DDS8888}-class single-phase energy IC: cumulative energy over Modbus with \textbf{0.01~kWh} (10~Wh) native register steps.

Firmware (\texttt{onepowerLS/onepwr-aws-mesh}) integrates active power (trapezoidal rule) to publish \texttt{EnergyIntegrated} at roughly \textbf{0.8~Wh} effective resolution alongside \texttt{EnergyActive}. A future PCB revision may route the DDS8888 \textbf{CF} pulse output to an ESP32 PCNT pin (1200~imp/kWh) to avoid integration drift; v2.2 has no hardware pulse input.

\section{Two gateway concepts (do not conflate)}

\subsection{Mesh uplink gateway (powerhouse, 1Meter fleet)}

At MAK, a device such as \textbf{23022667} is labelled ``gateway'' in ops SOPs: a \textbf{Mesh-Lite} infrastructure node at the powerhouse. It is \textbf{not} a customer meter and is not registered in \texttt{meters} for billing.

\textbf{Resilience:} mesh root is \textbf{not} a single named hardware choke point. Firmware semantics allow \textbf{any node with upstream Wi-Fi} to become root; the mesh can re-elect if the labelled gateway drops (\texttt{CONTEXT.md}).

\textbf{Role:} provide \textbf{IP reachability} to AWS (MQTT, OTA, telemetry). This path does \textbf{not} implement mobile-money settlement.

\subsection{Merchant payment gateway (new markets, no deep telecom integration)}

Entering a new market does \textbf{not} require a lengthy, costly integration with each mobile operator (Vodacom, MTN, bespoke USSD/short-code products, or operator billing APIs). The proven pattern in this program is:

\begin{enumerate}[nosep]
  \item Open a \textbf{standard merchant / agent mobile-money account} with the local provider.
  \item Install the \textbf{SMS gateway app} on the merchant Android phone (\texttt{SMS-Gateway-APP}): it forwards \textbf{transaction confirmation SMS} to the country's PHP \texttt{receive.php}.
  \item The PHP stack mirrors JSON to Customer Care; \textbf{\texttt{ingest.py}} parses confirmations $\rightarrow$ \textbf{1PDB} (\texttt{transactions}, etc.) $\rightarrow$ optional \textbf{SparkMeter credit}.
\end{enumerate}

This \textbf{closes the financial loop}---customer pays, confirmation lands on the merchant phone, structured ingest updates 1PDB---\textbf{without} coupling the minigrid platform to each carrier's telecom menu, USSD stack, or direct payment APIs. Lesotho and Benin production stacks are documented in \texttt{CONTEXT.md} (\texttt{SMSComms}, \texttt{SMSComms-BN}, \texttt{/api/sms/incoming}, \texttt{/api/bn/sms/incoming}).

\section{Connection architecture (meters and cloud)}

\begin{tabularx}{\linewidth}{@{}lX@{}}
\toprule
\textbf{Layer} & \textbf{Role} \\
\midrule
Edge & ESP32-C3, Mesh-Lite Wi-Fi; customer and infra nodes. \\
Cloud & AWS IoT Core (MQTT, TLS, OTA); \texttt{ingestion\_gate} Lambda pattern. \\
Integration & Lambda to CC \texttt{/api/meters/reading} with shared IoT key header. \\
Backend & 1PDB: \texttt{meter\_readings}, \texttt{hourly\_consumption} (\texttt{iot}), \texttt{prototype\_meter\_state}. \\
Phase 2+ & \texttt{relay\_control.py}: IoT relay command topic; Lambda forwards ack to CC; env gate \texttt{RELAY\_AUTO\_TRIGGER\_ENABLED}. \\
\bottomrule
\end{tabularx}

Devices report timestamps in \textbf{SAST (UTC+2)}; ingest stores \textbf{UTC}. Fleet offline monitoring may use DynamoDB \texttt{meter\_last\_seen} and \texttt{scripts/ops/monitor\_1meter\_offline.py}.

\section{1PDB and billing-source primacy}

1PDB is the system of record for customers, meters, payments, and consumption. For the 1Meter migration test, \texttt{balance\_engine.py} selects \textbf{one} consumption source per (account, hour): SparkMeter (\texttt{thundercloud} / \texttt{koios}) or 1Meter (\texttt{iot}), using \texttt{accounts.billing\_meter\_priority}, then \texttt{system\_config.billing\_meter\_priority} (fleet default \texttt{sm} in Phase~1). The non-primary source is shown as a what-if on Check Meters; it is not written to \texttt{transactions}.

\textbf{Hourly deltas in production ingest:} \texttt{ingest.py} uses \textbf{\texttt{energy\_active}} (Modbus cumulative) for deltas because it survives power cycles; \texttt{energy\_integrated} resets on ESP32 reboot and is not used for production hourly accumulation.

\section{Migration status (summary)}

\begin{description}[style=nextline]
  \item[Phase 1 (active)] SparkMeter primary billing; 1Meter check / what-if on MAK; fleet default \texttt{sm}.
  \item[Phase 2 (designed, gated)] Flip default to \texttt{1m}; SparkMeter guard-credit; CC-driven 1Meter relay with env gating.
  \item[Phase 3 (future)] 1Meter-only; retire SM check meters.
\end{description}

OTA and remote build tooling live on the documented firmware build host; full S3 release automation remains a next step per \texttt{CONTEXT.md}.

\section{Impact}

\textbf{Cost and ownership:} path to cheaper meters and no per-site vendor cloud lock-in for new rollouts. \textbf{Risk:} parallel SM+1M with explicit source priority avoids silent double-counting. \textbf{Markets:} merchant-gateway payment path reduces dependence on operator-specific integrations. \textbf{Operations:} mesh health, OTA discipline, and parse accuracy for SMS confirmations remain operational levers.

\section{References (in-repo)}

\noindent
\texttt{CONTEXT.md};

\noindent
\texttt{docs/ops/1meter-billing-migration-protocol.md};

\noindent
\texttt{docs/ops/1meter-build-and-ota-runbook.md}, \texttt{docs/ops/1meter-fleet-monitor.md};

\noindent
\texttt{acdb-api/ingest.py}, \texttt{balance\_engine.py}, \texttt{relay\_control.py}.

\vfill
\noindent\textit{Source file: \texttt{docs/ops/1meter-1pdb-white-paper.tex}. Regenerate PDF with \texttt{pdflatex}.}

\end{document}
